%!TEX program = pdflatex
% =====================================================================
%  VEKTORY – studijní text pro 1. ročník VUT FEKT (obor EKT)
%  Překlad: pdflatex (nebo Overleaf – kompilátor pdfLaTeX)
% =====================================================================
\documentclass[a4paper,11pt]{article}

% ---------- Kódování, jazyk, písmo ----------
\usepackage[T1]{fontenc}
\usepackage[utf8]{inputenc}
\usepackage[czech]{babel}
\usepackage{lmodern}
\usepackage[expansion=false]{microtype}

% ---------- Rozvržení stránky ----------
\usepackage[a4paper,top=2.6cm,bottom=2.6cm,left=2.5cm,right=2.5cm,headheight=24pt]{geometry}
\usepackage{setspace}
\setstretch{1.08}

% ---------- Matematika ----------
\usepackage{amsmath,amssymb,mathtools}

% ---------- Barvy ----------
\usepackage[dvipsnames,table]{xcolor}
\definecolor{primary}{HTML}{1F4E79}   % hlavní modrá
\definecolor{accent}{HTML}{B03A2E}    % výrazná červená (výsledné vektory)
\definecolor{veccol1}{HTML}{1F6FB2}   % modrá – první vektor
\definecolor{veccol2}{HTML}{2E8B57}   % zelená – druhý vektor
\definecolor{ektcol}{HTML}{C77700}    % oranžová – rámečky ECT

% ---------- Grafika ----------
\usepackage{tikz}
\usetikzlibrary{arrows.meta,calc,angles,babel}
\tikzset{
  vec/.style   ={-{Stealth[length=2.6mm,width=2mm]}, very thick, veccol1},
  vecg/.style  ={-{Stealth[length=2.6mm,width=2mm]}, very thick, veccol2},
  vecr/.style  ={-{Stealth[length=2.6mm,width=2mm]}, very thick, accent},
  ax/.style    ={-{Stealth[length=2.2mm,width=1.6mm]}, thick, gray!70!black},
  helpl/.style ={dashed, gray!70},
}

% ---------- Rámečky ----------
\usepackage[most]{tcolorbox}
\usepackage{etoolbox}

\newtcolorbox[auto counter,number within=section]{definice}[1]{%
  enhanced, sharp corners=south, arc=2mm,
  colback=primary!5, colframe=primary, boxrule=0.8pt,
  coltitle=white, colbacktitle=primary, fonttitle=\bfseries,
  title={Definice~\thetcbcounter\ \textnormal{---} #1},
  before skip=10pt, after skip=10pt}

\newtcolorbox[auto counter,number within=section]{veta}[1]{%
  enhanced, sharp corners=south, arc=2mm,
  colback=accent!5, colframe=accent!85!black, boxrule=0.8pt,
  coltitle=white, colbacktitle=accent!85!black, fonttitle=\bfseries,
  title={Věta~\thetcbcounter\ \textnormal{---} #1},
  before skip=10pt, after skip=10pt}

\newtcolorbox[auto counter,number within=section]{priklad}[1]{%
  enhanced, breakable, sharp corners=south, arc=2mm,
  colback=green!4!white, colframe=green!45!black, boxrule=0.8pt,
  coltitle=white, colbacktitle=green!45!black, fonttitle=\bfseries,
  title={Příklad~\thetcbcounter\ \textnormal{---} #1},
  before skip=10pt, after skip=10pt}

\newtcolorbox{poznamka}{%
  enhanced, frame hidden, sharp corners, boxrule=0pt,
  colback=gray!10, borderline west={3pt}{0pt}{gray!60},
  before upper={\textbf{Poznámka.}\ },
  before skip=8pt, after skip=8pt}

\newtcolorbox{pozor}{%
  enhanced, frame hidden, sharp corners, boxrule=0pt,
  colback=red!6, borderline west={3pt}{0pt}{red!70!black},
  before upper={\textbf{\color{red!70!black}Pozor!}\ },
  before skip=8pt, after skip=8pt}

\newtcolorbox{ekt}[1][Použití v elektrotechnice]{%
  enhanced, breakable, arc=2mm, boxrule=0.8pt,
  colback=ektcol!7, colframe=ektcol,
  coltitle=black, colbacktitle=ektcol!30, fonttitle=\bfseries\small,
  title={$\blacktriangleright$ #1},
  before skip=10pt, after skip=10pt}

% ---------- Seznamy, tabulky, odkazy ----------
\usepackage{float}
\usepackage{enumitem}
\setlist{itemsep=2pt,topsep=4pt}
\usepackage{booktabs}
\usepackage{array}
\renewcommand{\arraystretch}{1.25}

% ---------- Nadpisy ----------
\usepackage{titlesec}
\usepackage{needspace}
\titleformat{\section}
  {\Large\bfseries\color{primary}}{\thesection}{0.8em}{}[{\color{primary}\titlerule[1pt]}]
\titleformat{\subsection}
  {\large\bfseries\color{primary!85!black}}{\thesubsection}{0.7em}{}
\titleformat{\subsubsection}
  {\normalsize\bfseries\color{primary!70!black}}{\thesubsubsection}{0.6em}{}
\titlespacing*{\section}{0pt}{20pt}{10pt}
\titlespacing*{\subsection}{0pt}{14pt}{6pt}
\pretocmd{\section}{\Needspace*{16\baselineskip}}{}{}
\pretocmd{\subsection}{\Needspace*{6\baselineskip}}{}{}

% ---------- Záhlaví a zápatí ----------
\usepackage{fancyhdr}
\pagestyle{fancy}
\fancyhf{}
\fancyhead[L]{\small\color{primary}\textbf{Vektory}}
\fancyhead[R]{\small\color{primary}\nouppercase{\leftmark}}
\fancyfoot[C]{\small\thepage}
\renewcommand{\headrulewidth}{0.4pt}
\renewcommand{\headrule}{\hbox to\headwidth{\color{primary}\leaders\hrule height \headrulewidth\hfill}}

\usepackage[hypertexnames=false,pdfpagelabels,colorlinks=true,linkcolor=primary,urlcolor=primary,citecolor=primary,
            pdftitle={Vektory – studijní text},pdfauthor={Studijní materiál}]{hyperref}

% ---------- Vlastní příkazy ----------
\newcommand{\R}{\mathbb{R}}
\newcommand{\abs}[1]{\left\lvert #1\right\rvert}
\newcommand{\ii}{\vec{\imath}}
\newcommand{\jj}{\vec{\jmath}}
\newcommand{\kk}{\vec{k}}
\newcommand{\jed}[1]{\,\mathrm{#1}}      % jednotky: 5\jed{N}
\newcommand{\dc}{{,}}                    % desetinná čárka v matematice
\DeclareMathOperator{\proj}{proj}
\DeclareMathOperator{\atan}{arctg}
\newcommand{\jimag}{\mathrm{j}}          % imaginární jednotka (elektrotechnika)

\begin{document}

% =====================================================================
%  TITULNÍ STRANA
% =====================================================================
\begin{titlepage}
\centering
\vspace*{1.5cm}
{\color{primary}\rule{\textwidth}{2.5pt}}\\[0.8cm]
{\fontsize{46}{52}\selectfont\bfseries\color{primary} Vektory\par}
\vspace{0.4cm}
{\Large\color{primary!80!black} Základy vektorového počtu ve 2D a 3D\par}
\vspace{0.5cm}
{\color{primary}\rule{\textwidth}{2.5pt}}\\[1.2cm]

\begin{tikzpicture}[scale=1.15]
  \draw[ax] (-0.3,0) -- (6.3,0);
  \draw[ax] (0,-0.3) -- (0,3.6);
  \draw[vec]  (0,0) -- (4,1)   node[below right]{$\vec a$};
  \draw[vecg] (0,0) -- (1.5,3) node[above left]{$\vec b$};
  \draw[vecr] (0,0) -- (5.5,4) node[above right]{$\vec a+\vec b$};
  \draw[helpl] (4,1) -- (5.5,4);
  \draw[helpl] (1.5,3) -- (5.5,4);
\end{tikzpicture}

\vspace{1.4cm}
\begin{tcolorbox}[enhanced,colback=primary!5,colframe=primary!40,boxrule=0.6pt,arc=2mm,width=0.85\textwidth]
\centering\large
Studijní text pro 1.~ročník\\[2pt]
\textbf{Vysoké učení technické v Brně, FEKT}\\[2pt]
obor Elektronika a komunikační technologie (EKT)
\end{tcolorbox}

\vspace{1cm}
\begin{minipage}{0.8\textwidth}
\centering\normalsize
Co, jak a proč: pojem vektoru, souřadnice, sčítání, odčítání, násobení skalárem,
skalární, vektorový a smíšený součin, skládání a rozklad vektorů,
fázory. Všechno s obrázky, řešenými příklady a příklady k procvičení.
\end{minipage}

\vfill
{\large\today}
\end{titlepage}

% =====================================================================
%  OBSAH
% =====================================================================
\pagenumbering{roman}
\tableofcontents
\clearpage
\pagenumbering{arabic}

% =====================================================================
\section{Úvod a základní pojmy}
% =====================================================================

\subsection{Skalár versus vektor}

Fyzikální a technické veličiny se dělí na dvě velké skupiny podle toho, čím je
potřebujeme popsat.

\begin{itemize}
  \item \textbf{Skalární veličina (skalár)} je určena \emph{jediným číslem} (s~jednotkou).
        Příklady: teplota, hmotnost, čas, práce, energie, elektrický náboj, elektrický
        potenciál, výkon.
  \item \textbf{Vektorová veličina (vektor)} je určena \emph{velikostí a směrem} (včetně
        orientace). Příklady: rychlost, zrychlení, síla, intenzita elektrického pole
        $\vec E$, magnetická indukce $\vec B$, hustota proudu $\vec J$, moment síly.
\end{itemize}

Teplota v místnosti je $21\,^\circ\mathrm{C}$ a to je všechno. Ale o síle
$10\jed{N}$ nevíme nic, dokud neřekneme, \emph{kterým směrem} působí.

\begin{ekt}[Proč vektory v ECT]
Prakticky celá elektrotechnika stojí na vektorech: elektrické a magnetické pole
($\vec E,\ \vec D,\ \vec B,\ \vec H$), Maxwellovy rovnice, Lorentzova síla,
Poyntingův vektor (tok energie elektromagnetické vlny), střídavé obvody
(fázory), zpracování signálů (signál jako vektor v prostoru vysoké dimenze),
komunikační technologie (konstelační diagramy, IQ modulace). Vektorový počet je
proto jedním ze základních nástrojů, které budete v dalších semestrech
používat neustále.
\end{ekt}

\subsection{Geometrický vektor}

\begin{definice}{Vektor}
\textbf{Orientovaná úsečka} $\overrightarrow{AB}$ je úsečka s~počátečním bodem $A$
a koncovým bodem $B$. \textbf{Vektor} $\vec a$ je množina všech orientovaných úseček,
které mají stejnou \emph{délku}, stejný \emph{směr} (jsou rovnoběžné) a stejnou
\emph{orientaci}. Každá z~těchto úseček je \emph{umístěním} vektoru $\vec a$.
\end{definice}

Vektor tedy není „připevněn“ k~žádnému konkrétnímu bodu: můžeme ho libovolně
posouvat (rovnoběžně), aniž by se změnil. Takový vektor se nazývá \emph{volný}.
Naopak u~některých úloh (síla působící v~konkrétním bodě) záleží na \emph{působišti}
a mluvíme o~vázaném vektoru.

\begin{figure}[!htbp]
\centering
\begin{tikzpicture}[scale=1.0]
  \coordinate (A) at (0,0);
  \coordinate (B) at (4,1.5);
  \draw[vec] (A) -- (B) node[midway,above left]{$\vec a=\overrightarrow{AB}$};
  \fill (A) circle (2pt) node[below left]{$A$ \small(počátek)};
  \fill (B) circle (2pt) node[above right]{$B$ \small(konec)};
  % posunuté umístění
  \draw[vec,opacity=0.55] (1,-1.6) -- (5,-0.1) node[midway,below right]{$\vec a$};
  \draw[helpl] (A) -- (1,-1.6);
  \draw[helpl] (B) -- (5,-0.1);
\end{tikzpicture}
\caption{Vektor $\vec a$ je určen délkou, směrem a orientací. Posunutím vznikne jiné
umístění téhož vektoru.}
\end{figure}

\subsection{Základní pojmy}

\begin{description}[leftmargin=1.2em,style=nextline,font=\bfseries\color{primary}]
  \item[Velikost (délka, norma) vektoru] Značí se $\abs{\vec a}$ nebo $a$.
        Je to vždy nezáporné číslo.
  \item[Nulový vektor $\vec o$] Má velikost $0$, nemá určitý směr. Je to jediný vektor,
        který je „bez směru“.
  \item[Jednotkový vektor] Vektor s~velikostí $1$. Jednotkový vektor ve směru $\vec a$
        značíme $\vec a^{\,0}$ a získáme jej \emph{normováním}:
        \[
          \vec a^{\,0}=\frac{\vec a}{\abs{\vec a}},\qquad \vec a\neq\vec o .
        \]
  \item[Opačný vektor $-\vec a$] Stejná velikost i~směr, opačná orientace.
  \item[Rovnost vektorů] $\vec a=\vec b$ právě tehdy, když mají stejnou velikost, směr a
        orientaci (tj. jsou to dvě umístění téhož vektoru).
  \item[Kolineární vektory] Leží na jedné přímce nebo na rovnoběžných přímkách
        (jsou rovnoběžné). Nulový vektor je kolineární s~každým.
  \item[Komplanární vektory] Leží v~jedné rovině nebo v~rovnoběžných rovinách
        (tři a více vektorů v~prostoru).
  \item[Polohový vektor bodu] Vektor $\vec r_P=\overrightarrow{OP}$ z~počátku
        souřadnic $O$ do bodu $P$. Polohu bodu popisuje jednoznačně.
\end{description}

\begin{pozor}
Vektor a~jeho velikost jsou \emph{různé objekty}. Vektor $\vec a$ není číslo, a proto
nelze psát například $\vec a>\vec b$ ani „dělit“ vektorem. Číslo je $\abs{\vec a}$.
\end{pozor}

\begin{poznamka}
Zápis vektorů: v~rukopise se používá šipka $\vec a$ (nebo podtržení či vlnovka),
v~tištěných textech tučné písmo $\mathbf a$. V~tomto textu používáme šipku.
Vektor mezi body $A$ a $B$ zapisujeme $\overrightarrow{AB}$.
\end{poznamka}

% =====================================================================
\section{Vektory v~rovině (2D)}
% =====================================================================

\subsection{Kartézská soustava souřadnic a báze}

V~rovině zvolíme dvě navzájem kolmé osy $x$ a $y$ se společným počátkem $O$ a na
nich \emph{jednotkové vektory} $\ii$ (ve směru osy $x$) a $\jj$ (ve směru osy $y$).
Dvojice $(\ii,\jj)$ se nazývá \textbf{ortonormální báze} (kolmé vektory o~velikosti
$1$). Každý vektor v~rovině lze jednoznačně zapsat jako
\begin{equation}
  \vec a = a_x\,\ii + a_y\,\jj = (a_x,\,a_y),
\end{equation}
kde čísla $a_x,a_y$ jsou \textbf{souřadnice} (složky) vektoru. Vektory
$a_x\ii$ a~$a_y\jj$ jsou \emph{vektorové složky} a čísla $a_x,a_y$ jsou
\emph{skalární složky} (průměty na osy).

\begin{figure}[!htbp]
\centering
\begin{tikzpicture}[scale=1.15]
  \coordinate (O) at (0,0);
  \coordinate (X) at (4,0);
  \coordinate (P) at (4,3);
  \draw[ax] (-0.5,0) -- (5.3,0) node[right]{$x$};
  \draw[ax] (0,-0.5) -- (0,4) node[above]{$y$};
  \draw[helpl] (0,3) -- (P);
  \draw[helpl] (X) -- (P);
  \draw[vecg,opacity=0.8] (O) -- (X) node[midway,below]{$a_x\,\ii$};
  \draw[vecg,opacity=0.8] (X) -- (P) node[midway,right]{$a_y\,\jj$};
  \draw[vecr] (O) -- (P) node[midway,above left]{$\vec a$};
  \draw[vec] (O) -- (1,0) node[below]{$\ii$};
  \draw[vec] (O) -- (0,1) node[left]{$\jj$};
  \draw pic[draw,pic text={\(\varphi\)},angle radius=9mm,angle eccentricity=1.4] {angle=X--O--P};
  \node[left] at (0,3) {$a_y$};
  \node[below] at (4,-0.4) {$a_x$};
\end{tikzpicture}
\caption{Vektor $\vec a=(a_x,a_y)$ v~rovině a jeho rozklad do báze $(\ii,\jj)$.}
\end{figure}

\subsection{Velikost a jednotkový vektor}

Z~Pythagorovy věty (pravoúhlý trojúhelník se stranami $a_x,a_y$ a přeponou
$\abs{\vec a}$) plyne
\begin{equation}
  \abs{\vec a}=\sqrt{a_x^2+a_y^2},\qquad
  \vec a^{\,0}=\frac{1}{\abs{\vec a}}\,(a_x,a_y).
\end{equation}

\subsection{Vektor určený dvěma body}

Jsou-li dány body $A=[x_A,y_A]$ a $B=[x_B,y_B]$, pak
\begin{equation}
  \overrightarrow{AB}=B-A=(x_B-x_A,\;y_B-y_A),\qquad
  \abs{\overrightarrow{AB}}=\sqrt{(x_B-x_A)^2+(y_B-y_A)^2}.
\end{equation}
Pamatujte: \emph{konec minus začátek}.

\subsection{Směr vektoru a polární tvar}

Směr vektoru v~rovině udává úhel $\varphi$ od kladné poloosy $x$ (kladně proti směru
hodinových ručiček). Platí
\begin{equation}
  a_x=\abs{\vec a}\cos\varphi,\qquad a_y=\abs{\vec a}\sin\varphi,\qquad
  \operatorname{tg}\varphi=\frac{a_y}{a_x}\ \ (a_x\neq0).
\end{equation}

\begin{pozor}
Funkce $\operatorname{arctg}$ na kalkulačce vrací úhel jen v~intervalu
$(-90^\circ,90^\circ)$. Kvadrant určujte podle znamének $a_x$ a $a_y$! Pro
$a_x<0$ je třeba přičíst $180^\circ$ ($\pi$). Řada kalkulaček a programů má pro
to funkci \texttt{atan2(y, x)} nebo převod \texttt{Pol(x, y)}.
\end{pozor}

\begin{ekt}[Polární tvar a komplexní čísla]
Vektor $(a_x,a_y)$ v~rovině je totéž co komplexní číslo $a_x+\jimag\,a_y$.
Převod mezi kartézským a polárním tvarem, který jste právě viděli, budete
používat při počítání s~fázory ve střídavých obvodech (kapitola~\ref{sec:skladani}).
\end{ekt}

\begin{priklad}{Velikost a směr vektoru}
Určete velikost a směr vektoru $\vec a=(-3,\,3)$ a jednotkový vektor $\vec a^{\,0}$.

\medskip\noindent\textbf{Řešení.}
\[
  \abs{\vec a}=\sqrt{(-3)^2+3^2}=\sqrt{18}=3\sqrt2\doteq4\dc24 .
\]
Je $\operatorname{tg}\varphi=3/(-3)=-1$, ale vektor leží ve druhém kvadrantu
($a_x<0,\ a_y>0$), proto $\varphi=180^\circ-45^\circ=135^\circ$.
Jednotkový vektor:
\[
  \vec a^{\,0}=\frac{(-3,3)}{3\sqrt2}=\left(-\frac{1}{\sqrt2},\frac{1}{\sqrt2}\right)
  \doteq(-0\dc707,\;0\dc707).
\]
\end{priklad}

\begin{priklad}{Vektor mezi dvěma body}
Body $A=[2,-1]$, $B=[5,3]$. Určete $\overrightarrow{AB}$ a jeho velikost.

\medskip\noindent\textbf{Řešení.}
$\overrightarrow{AB}=(5-2,\;3-(-1))=(3,4)$, $\abs{\overrightarrow{AB}}=\sqrt{9+16}=5$.
\end{priklad}

\subsection{Kolmý vektor v~rovině}

Otočením vektoru $\vec a=(a_x,a_y)$ o~$90^\circ$ proti směru hodinových ručiček
vznikne vektor $(-a_y,\,a_x)$, o~$90^\circ$ po směru $(a_y,-a_x)$. Oba jsou kolmé na
$\vec a$ a mají stejnou velikost. Trik je užitečný například při hledání normálového
vektoru přímky.

% =====================================================================
\section{Vektory v~prostoru (3D)}
% =====================================================================

\subsection{Pravotočivá soustava souřadnic}

V~prostoru zavedeme tři navzájem kolmé osy $x,y,z$ a jednotkové vektory
$\ii,\jj,\kk$. Soustava je \textbf{pravotočivá}, jestliže platí \emph{pravidlo pravé ruky}:
míří-li palec ve směru osy $x$ a ukazováček ve směru osy $y$, míří prostředníček
(kolmo na obě) ve směru osy $z$. Kdybychom se dívali z~kladné poloosy $z$ na
rovinu $xy$, otočení z~osy $x$ do osy $y$ (o $90^\circ$) je proti směru hodinových
ručiček. V~celém textu se používá pravotočivá soustava.

\begin{figure}[!htbp]
\centering
\begin{tikzpicture}[x={(-0.55cm,-0.42cm)},y={(1.1cm,0cm)},z={(0cm,1.1cm)}]
  \coordinate (O) at (0,0,0);
  \coordinate (P) at (3,3,2.5);
  % osy
  \draw[ax] (0,0,0) -- (5,0,0) node[below left]{$x$};
  \draw[ax] (0,0,0) -- (0,4.6,0) node[right]{$y$};
  \draw[ax] (0,0,0) -- (0,0,3.9) node[above]{$z$};
  % kvádr
  \draw[helpl] (3,0,0) -- (3,3,0) -- (0,3,0);
  \draw[helpl] (3,3,0) -- (P);
  \draw[helpl] (0,0,2.5) -- (0,3,2.5) -- (P);
  \draw[helpl] (3,0,0) -- (3,0,2.5) -- (P);
  \draw[helpl] (0,3,0) -- (0,3,2.5);
  \draw[helpl] (0,0,2.5) -- (3,0,2.5);
  % složky
  \draw[vecg,opacity=0.85] (0,0,0) -- (3,0,0) node[midway,below left]{$a_x\ii$};
  \draw[vecg,opacity=0.85] (0,0,0) -- (0,3,0) node[midway,below right]{$a_y\jj$};
  \draw[vecg,opacity=0.85] (0,0,0) -- (0,0,2.5) node[midway,left]{$a_z\kk$};
  % vektor
  \draw[vecr] (O) -- (P) node[above right]{$\vec a=(a_x,a_y,a_z)$};
  \fill (P) circle (1.6pt);
\end{tikzpicture}
\caption{Pravotočivá soustava a vektor $\vec a$ v~prostoru. Osa $x$ míří „k~nám“.}
\end{figure}

\subsection{Souřadnice, velikost, vzdálenost}

Každý vektor v~prostoru má jednoznačný zápis
\begin{equation}
  \vec a=a_x\,\ii+a_y\,\jj+a_z\,\kk=(a_x,a_y,a_z).
\end{equation}
Dvojím užitím Pythagorovy věty dostaneme
\begin{equation}
  \abs{\vec a}=\sqrt{a_x^2+a_y^2+a_z^2}.
\end{equation}

Pro body $A=[x_A,y_A,z_A]$, $B=[x_B,y_B,z_B]$ platí
\begin{align}
  \overrightarrow{AB}&=(x_B-x_A,\;y_B-y_A,\;z_B-z_A),\\
  \abs{\overrightarrow{AB}}&=\sqrt{(x_B-x_A)^2+(y_B-y_A)^2+(z_B-z_A)^2},\\
  S_{AB}&=\left[\frac{x_A+x_B}{2},\ \frac{y_A+y_B}{2},\ \frac{z_A+z_B}{2}\right]
  \quad\text{(střed úsečky).}
\end{align}

\subsection{Směrové kosiny}

Označme $\alpha,\beta,\gamma$ úhly, které vektor $\vec a\neq\vec o$ svírá po řadě
s~kladnými poloosami $x,y,z$. Čísla $\cos\alpha,\cos\beta,\cos\gamma$ se nazývají
\textbf{směrové kosiny}:
\begin{equation}
  \cos\alpha=\frac{a_x}{\abs{\vec a}},\quad
  \cos\beta=\frac{a_y}{\abs{\vec a}},\quad
  \cos\gamma=\frac{a_z}{\abs{\vec a}},\qquad
  \cos^2\alpha+\cos^2\beta+\cos^2\gamma=1 .
\end{equation}
Jsou to zároveň souřadnice jednotkového vektoru:
$\vec a^{\,0}=(\cos\alpha,\cos\beta,\cos\gamma)$.

\begin{priklad}{Vektor v prostoru}
Jsou dány body $A=[1,2,-1]$ a $B=[4,-2,3]$. Určete $\overrightarrow{AB}$, jeho velikost,
směrové kosiny a střed úsečky $AB$.

\medskip\noindent\textbf{Řešení.}
\begin{align*}
  \overrightarrow{AB}&=(4-1,\,-2-2,\,3-(-1))=(3,-4,4),\\
  \abs{\overrightarrow{AB}}&=\sqrt{9+16+16}=\sqrt{41}\doteq6\dc40,\\
  \cos\alpha&=\tfrac{3}{\sqrt{41}}\doteq0\dc468\ \Rightarrow\ \alpha\doteq62\dc1^\circ,\\
  \cos\beta&=\tfrac{-4}{\sqrt{41}}\doteq-0\dc625\ \Rightarrow\ \beta\doteq128\dc7^\circ,\\
  \cos\gamma&=\tfrac{4}{\sqrt{41}}\doteq0\dc625\ \Rightarrow\ \gamma\doteq51\dc3^\circ,\\
  S&=\left[\tfrac{1+4}{2},\tfrac{2-2}{2},\tfrac{-1+3}{2}\right]=[2\dc5,\,0,\,1].
\end{align*}
Kontrola: $\tfrac{9}{41}+\tfrac{16}{41}+\tfrac{16}{41}=1$. \checkmark
\end{priklad}

\subsection{Zobecnění do \texorpdfstring{$\R^n$}{R\textasciicircum n}}

Všechno, co jsme řekli, lze bez změny zobecnit na uspořádané $n$-tice
$\vec a=(a_1,a_2,\dots,a_n)\in\R^n$, s~velikostí
$\abs{\vec a}=\sqrt{a_1^2+\dots+a_n^2}$. Ve 2D a 3D lze vektory kreslit, ve vyšších
dimenzích je popisujeme jen algebraicky. V~teorii signálů či strojovém učení
pracujete s~vektory o~tisících složek, ale počítá se s~nimi \emph{přesně stejně}.
Výjimkou je \emph{vektorový součin}, který je (v~této podobě) specifický pro $\R^3$.

% =====================================================================
\section{Sčítání a odčítání vektorů}
% =====================================================================

\subsection{Geometrický význam}

\begin{definice}{Součet vektorů}
Součet $\vec a+\vec b$ získáme jedním ze dvou ekvivalentních způsobů:
\begin{itemize}
  \item \textbf{Trojúhelníkové pravidlo:} vektor $\vec b$ umístíme tak, aby začínal
        v~koncovém bodě vektoru $\vec a$. Součet $\vec a+\vec b$ vede z~počátku
        vektoru $\vec a$ do konce vektoru $\vec b$.
  \item \textbf{Rovnoběžníkové pravidlo:} oba vektory umístíme do společného počátku.
        Součet je úhlopříčka rovnoběžníku, který vektory určují, vycházející
        z~téhož počátku.
\end{itemize}
\end{definice}

\begin{figure}[!htbp]
\centering
\begin{tikzpicture}[scale=1.0]
  % Trojúhelníkové pravidlo
  \begin{scope}
    \draw[vec]  (0,0) -- (3,0.5) node[midway,below]{$\vec a$};
    \draw[vecg] (3,0.5) -- (4.5,2.5) node[midway,right]{$\vec b$};
    \draw[vecr] (0,0) -- (4.5,2.5) node[midway,above left]{$\vec a+\vec b$};
    \node at (2.2,-1.3) {\small trojúhelníkové pravidlo};
  \end{scope}
  % Rovnoběžníkové pravidlo
  \begin{scope}[shift={(7.2,0)}]
    \draw[vec]  (0,0) -- (3,0.5) node[midway,below]{$\vec a$};
    \draw[vecg] (0,0) -- (1.5,2) node[midway,left]{$\vec b$};
    \draw[helpl] (3,0.5) -- (4.5,2.5);
    \draw[helpl] (1.5,2) -- (4.5,2.5);
    \draw[vecr] (0,0) -- (4.5,2.5) node[near end,above left]{$\vec a+\vec b$};
    \node at (2.2,-1.3) {\small rovnoběžníkové pravidlo};
  \end{scope}
\end{tikzpicture}
\caption{Dvě ekvivalentní konstrukce součtu vektorů.}
\end{figure}

Trojúhelníkové pravidlo se snadno rozšíří na \emph{součet více vektorů}
(mnohoúhelníkové pravidlo): vektory řadíme „konec na začátek“ a součet vede
z~počátku prvního do konce posledního. Pořadí přitom nehraje roli.

\subsection{Odčítání}

Rozdíl definujeme jako přičtení opačného vektoru: $\vec a-\vec b=\vec a+(-\vec b)$.
Když oba vektory umístíme do společného počátku, je \textbf{rozdíl $\vec a-\vec b$
vektor mířící z~konce vektoru $\vec b$ do konce vektoru $\vec a$}. (Proto také platí
$\overrightarrow{AB}=\vec r_B-\vec r_A$.)

\begin{figure}[!htbp]
\centering
\begin{tikzpicture}[scale=1.1]
  \draw[vec]  (0,0) -- (4,2.2) node[midway,above left]{$\vec a$};
  \draw[vecg] (0,0) -- (1.5,2.8) node[midway,left]{$\vec b$};
  \draw[vecr] (1.5,2.8) -- (4,2.2) node[midway,above]{$\vec a-\vec b$};
  \draw[vecg,dashed,opacity=0.6] (4,2.2) -- (2.5,-0.6) node[midway,right]{$-\vec b$};
  \draw[vecr,dashed,opacity=0.6] (0,0) -- (2.5,-0.6);
\end{tikzpicture}
\caption{Rozdíl vektorů: $\vec a-\vec b$ míří od konce $\vec b$ ke konci $\vec a$.
Přerušovaně: konstrukce jako $\vec a+(-\vec b)$.}
\end{figure}

\subsection{Výpočet pomocí souřadnic}

Vektory se sčítají a odčítají \textbf{po složkách}:
\begin{align}
  \vec a\pm\vec b&=(a_x\pm b_x,\;a_y\pm b_y)&&\text{(2D)},\\
  \vec a\pm\vec b&=(a_x\pm b_x,\;a_y\pm b_y,\;a_z\pm b_z)&&\text{(3D)}.
\end{align}

\begin{veta}{Vlastnosti sčítání vektorů}
Pro libovolné vektory $\vec a,\vec b,\vec c$ platí:
\begin{enumerate}[label=(\roman*),nosep]
  \item $\vec a+\vec b=\vec b+\vec a$ \hfill (komutativita)
  \item $(\vec a+\vec b)+\vec c=\vec a+(\vec b+\vec c)$ \hfill (asociativita)
  \item $\vec a+\vec o=\vec a$ \hfill (nulový vektor)
  \item $\vec a+(-\vec a)=\vec o$ \hfill (opačný vektor)
\end{enumerate}
\end{veta}

\begin{veta}{Trojúhelníková nerovnost}
$\abs{\,\abs{\vec a}-\abs{\vec b}\,}\ \le\ \abs{\vec a\pm\vec b}\ \le\ \abs{\vec a}+\abs{\vec b}.$
Rovnost vpravo nastane, jsou-li vektory \emph{souhlasně} rovnoběžné.
\end{veta}

\begin{pozor}
Velikost součtu \emph{není} součet velikostí! Například $\vec a=(3,0)$,
$\vec b=(0,4)$: $\abs{\vec a+\vec b}=5$, ale $\abs{\vec a}+\abs{\vec b}=7$.
\end{pozor}

\begin{priklad}{Sčítání a odčítání ve 2D a 3D}
\textbf{a)} $\vec a=(2,5)$, $\vec b=(-1,3)$:\quad
$\vec a+\vec b=(1,8)$, \quad $\vec a-\vec b=(3,2)$, \quad
$2\vec a-3\vec b=(4+3,\,10-9)=(7,1)$.

\smallskip\noindent
\textbf{b)} $\vec u=(2,-1,3)$, $\vec v=(1,4,-2)$:\quad
$\vec u+\vec v=(3,3,1)$, \quad $\vec u-\vec v=(1,-5,5)$,\quad
$3\vec u-2\vec v=(6-2,\,-3-8,\,9+4)=(4,-11,13)$.
\end{priklad}

% =====================================================================
\section{Násobení vektoru skalárem}
% =====================================================================

\begin{definice}{Násobek vektoru}
Součin čísla $k\in\R$ a vektoru $\vec a$ je vektor $k\vec a$, který má
\begin{itemize}
  \item velikost $\abs{k\vec a}=\abs{k}\,\abs{\vec a}$,
  \item stejný směr jako $\vec a$, orientaci \emph{stejnou} pro $k>0$ a \emph{opačnou}
        pro $k<0$ (pro $k=0$ je výsledek nulový vektor).
\end{itemize}
Po složkách: $k\vec a=(ka_x,\,ka_y,\,ka_z)$.
\end{definice}

\begin{figure}[!htbp]
\centering
\begin{tikzpicture}[scale=1.15]
  \draw[vec]  (0,3.2) -- (1.5,3.95) node[right]{$\vec a$};
  \draw[vec]  (0,2.0) -- (3,3.5)    node[right]{$2\vec a$};
  \draw[vec]  (0,0.9) -- (0.75,1.275) node[right]{$\tfrac12\vec a$};
  \draw[vecr] (4.5,0) -- (3,-0.75)  node[left]{$-\vec a$};
\end{tikzpicture}
\caption{Násobky vektoru $\vec a$: změna velikosti a případně orientace, směr zůstává.}
\end{figure}

\begin{veta}{Vlastnosti násobení skalárem}
Pro $k,l\in\R$ a vektory $\vec a,\vec b$ platí
\[
  k(l\vec a)=(kl)\vec a,\quad
  (k+l)\vec a=k\vec a+l\vec a,\quad
  k(\vec a+\vec b)=k\vec a+k\vec b,\quad
  1\cdot\vec a=\vec a .
\]
\end{veta}

\begin{poznamka}
Vektory se \textbf{nedělí}. Výraz $\vec a/k$ ovšem smysl má: je to
$\tfrac1k\vec a$. Dělení vektoru vektorem neexistuje.
\end{poznamka}

\subsection{Vektor dané velikosti v~daném směru}

Chceme-li vektor o~velikosti $c$ ve směru $\vec a$, vezmeme jednotkový vektor
a natáhneme jej:
\begin{equation}
  \vec b=c\,\vec a^{\,0}=c\,\frac{\vec a}{\abs{\vec a}} .
\end{equation}

\begin{priklad}{Síla daného směru}
Síla o~velikosti $35\jed{N}$ působí ve směru vektoru $\vec s=(2,-3,6)$. Určete její
souřadnice.

\medskip\noindent\textbf{Řešení.} $\abs{\vec s}=\sqrt{4+9+36}=7$, tedy
\[
  \vec F=35\cdot\frac{(2,-3,6)}{7}=5\cdot(2,-3,6)=(10,\,-15,\,30)\jed{N}.
\]
\end{priklad}

\subsection{Kolinearita}

Nenulové vektory $\vec a,\vec b$ jsou \textbf{kolineární} (rovnoběžné) právě tehdy,
když existuje $k\in\R$ s~$\vec b=k\vec a$; ekvivalentně, když jsou jejich souřadnice
\emph{úměrné}:
\[
  \frac{b_x}{a_x}=\frac{b_y}{a_y}=\frac{b_z}{a_z}\;(=k).
\]

\begin{priklad}{Kolineární vektory}
$\vec a=(2,-4,6)$, $\vec b=(-1,2,-3)$.
Poměry: $\tfrac{-1}{2}=\tfrac{2}{-4}=\tfrac{-3}{6}=-\tfrac12$, tedy $\vec b=-\tfrac12\vec a$.
Vektory jsou kolineární, mají opačnou orientaci.
\end{priklad}

\subsection{Lineární kombinace, lineární závislost a báze}

\textbf{Lineární kombinace} vektorů $\vec a_1,\dots,\vec a_n$ je vektor
$k_1\vec a_1+\dots+k_n\vec a_n$ s~reálnými koeficienty $k_i$. Vektory jsou
\textbf{lineárně nezávislé}, pokud žádný z~nich nelze vyjádřit jako lineární kombinaci
ostatních; tj. rovnost $k_1\vec a_1+\dots+k_n\vec a_n=\vec o$ platí jen pro
$k_1=\dots=k_n=0$. V~opačném případě jsou \textbf{lineárně závislé}.

\begin{itemize}
  \item Dva vektory jsou lineárně závislé $\iff$ jsou \emph{kolineární}.
  \item Tři vektory v~prostoru jsou lineárně závislé $\iff$ jsou \emph{komplanární}
        (viz smíšený součin, sekce~\ref{sec:smiseny}).
  \item V~rovině tvoří \textbf{bázi} libovolné dva nekolineární vektory, v~prostoru
        libovolné tři nekomplanární vektory. Každý vektor pak lze \emph{jediným
        způsobem} vyjádřit jako jejich lineární kombinaci.
\end{itemize}

\begin{priklad}{Rozklad vektoru do jiné báze}
Vyjádřete $\vec c=(4,1)$ jako lineární kombinaci vektorů $\vec u=(1,1)$ a $\vec v=(1,-1)$.

\medskip\noindent\textbf{Řešení.} Hledáme $\alpha,\beta$ tak, aby
$\vec c=\alpha\vec u+\beta\vec v$. Po složkách
\[
  \begin{aligned}\alpha+\beta&=4\\ \alpha-\beta&=1\end{aligned}
  \quad\Longrightarrow\quad \alpha=\tfrac52,\ \beta=\tfrac32 .
\]
Tedy $\vec c=\tfrac52\vec u+\tfrac32\vec v$. Kontrola: $\tfrac52(1,1)+\tfrac32(1,-1)=(4,1)$. \checkmark
\end{priklad}

% =====================================================================
\section{Skalární součin}
% =====================================================================

\subsection{Definice a geometrický význam}

\begin{definice}{Skalární součin}
Nechť $\varphi\in\langle0^\circ,180^\circ\rangle$ je úhel, který svírají nenulové vektory
$\vec a,\vec b$. Jejich \textbf{skalární součin} je \emph{číslo}
\begin{equation}
  \vec a\cdot\vec b=\abs{\vec a}\,\abs{\vec b}\cos\varphi .
\end{equation}
Je-li aspoň jeden z~vektorů nulový, je $\vec a\cdot\vec b=0$.
\end{definice}

Výsledkem je \textbf{skalár}, ne vektor. Skalární součin měří, „jak moc míří vektory
stejným směrem“.

\begin{figure}[!htbp]
\centering
\begin{tikzpicture}[scale=1.15]
  \coordinate (O) at (0,0);
  \coordinate (B) at (5,0);
  \coordinate (A) at (3.4,2.4);
  \coordinate (Ap) at (3.4,0);
  \draw[vecg] (O) -- (B) node[right]{$\vec b$};
  \draw[vec]  (O) -- (A) node[above right]{$\vec a$};
  \draw[helpl] (A) -- (Ap);
  \draw[ultra thick,accent] (O) -- (Ap) node[midway,below=3pt]{$\abs{\vec a}\cos\varphi$};
  \draw (3.1,0) -- (3.1,0.3) -- (3.4,0.3);
  \draw pic[draw,pic text={\(\varphi\)},angle radius=9mm,angle eccentricity=1.4] {angle=B--O--A};
\end{tikzpicture}
\caption{$\vec a\cdot\vec b=\abs{\vec b}\cdot(\abs{\vec a}\cos\varphi)$ = velikost $\vec b$ krát
průmět $\vec a$ do směru $\vec b$.}
\end{figure}

\subsection{Výpočet pomocí souřadnic}

Protože $\ii\cdot\ii=\jj\cdot\jj=\kk\cdot\kk=1$ (úhel $0^\circ$) a
$\ii\cdot\jj=\jj\cdot\kk=\kk\cdot\ii=0$ (úhel $90^\circ$), po roznásobení
$(a_x\ii+a_y\jj+a_z\kk)\cdot(b_x\ii+b_y\jj+b_z\kk)$ zůstanou jen „stejné“ členy:
\begin{equation}
  \boxed{\ \vec a\cdot\vec b=a_xb_x+a_yb_y+a_zb_z\ }
  \qquad(\text{ve 2D bez třetího členu}).
\end{equation}

\subsection{Vlastnosti}

\begin{veta}{Vlastnosti skalárního součinu}
Pro libovolné vektory $\vec a,\vec b,\vec c$ a číslo $k$:
\begin{enumerate}[label=(\roman*),nosep]
  \item $\vec a\cdot\vec b=\vec b\cdot\vec a$ \hfill (komutativita)
  \item $\vec a\cdot(\vec b+\vec c)=\vec a\cdot\vec b+\vec a\cdot\vec c$ \hfill (distributivita)
  \item $(k\vec a)\cdot\vec b=k\,(\vec a\cdot\vec b)=\vec a\cdot(k\vec b)$
  \item $\vec a\cdot\vec a=\abs{\vec a}^2$, tedy $\abs{\vec a}=\sqrt{\vec a\cdot\vec a}$
  \item $\vec a\cdot\vec a\ge0$ a $\vec a\cdot\vec a=0\iff\vec a=\vec o$
  \item \textbf{Cauchyho--Schwarzova nerovnost:}\ $\abs{\vec a\cdot\vec b}\le\abs{\vec a}\abs{\vec b}$
\end{enumerate}
\end{veta}

\begin{pozor}
Skalární součin \emph{není asociativní} (výraz $(\vec a\cdot\vec b)\cdot\vec c$ nemá smysl:
skalár krát vektor dá vektor ve směru $\vec c$, ne $\vec a$). Ze vztahu
$\vec a\cdot\vec b=\vec a\cdot\vec c$ také \emph{nelze} krátit $\vec a$, tedy nevyplývá
$\vec b=\vec c$ (stačí, aby se $\vec b,\vec c$ lišily o~vektor kolmý na $\vec a$).
\end{pozor}

\subsection{Úhel vektorů a kolmost}

Z~definice vyjádříme kosinus úhlu:
\begin{equation}
  \cos\varphi=\frac{\vec a\cdot\vec b}{\abs{\vec a}\,\abs{\vec b}}
  =\frac{a_xb_x+a_yb_y+a_zb_z}{\sqrt{a_x^2+a_y^2+a_z^2}\ \sqrt{b_x^2+b_y^2+b_z^2}} .
\end{equation}
Znaménko skalárního součinu prozradí typ úhlu:
\[
  \vec a\cdot\vec b>0\ \Leftrightarrow\ \varphi\text{ ostrý},\qquad
  \vec a\cdot\vec b=0\ \Leftrightarrow\ \varphi=90^\circ,\qquad
  \vec a\cdot\vec b<0\ \Leftrightarrow\ \varphi\text{ tupý}.
\]

\begin{veta}{Podmínka kolmosti}
Nenulové vektory $\vec a,\vec b$ jsou \textbf{kolmé} ($\vec a\perp\vec b$) právě tehdy, když
$\vec a\cdot\vec b=0$.
\end{veta}

\subsection{Průmět vektoru do směru jiného vektoru}

Skalární průmět (se znaménkem) vektoru $\vec a$ do směru $\vec b$ a odpovídající
vektorová složka jsou
\begin{equation}
  a_{\parallel b}=\frac{\vec a\cdot\vec b}{\abs{\vec b}},\qquad
  \proj_{\vec b}\vec a=\frac{\vec a\cdot\vec b}{\vec b\cdot\vec b}\,\vec b
  =\frac{\vec a\cdot\vec b}{\abs{\vec b}^{2}}\,\vec b .
\end{equation}
Zbytek $\vec a_\perp=\vec a-\proj_{\vec b}\vec a$ je kolmý na $\vec b$.
Vektor $\vec a$ se tak rozpadá na složku rovnoběžnou a složku kolmou k~$\vec b$:
$\vec a=\proj_{\vec b}\vec a+\vec a_\perp$.

\begin{ekt}[Práce, výkon, tok]
\begin{itemize}
  \item \textbf{Práce} síly $\vec F$ po dráze $\vec s$: $W=\vec F\cdot\vec s=Fs\cos\varphi$.
  \item \textbf{Okamžitý výkon}: $P=\vec F\cdot\vec v$; v~elektrotechnice
        $p=\vec J\cdot\vec E$ (objemová hustota ztrát).
  \item \textbf{Napětí v~homogenním poli}: $U=\vec E\cdot\vec l$.
  \item \textbf{Magnetický tok plochou}: $\Phi=\vec B\cdot\vec S=BS\cos\varphi$, kde $\vec S$ je
        vektor plochy (velikost = obsah, směr = normála).
\end{itemize}
\end{ekt}

\begin{priklad}{Úhel mezi vektory}
$\vec a=(1,2,2)$, $\vec b=(2,-1,2)$.

\medskip\noindent\textbf{Řešení.}
$\vec a\cdot\vec b=2-2+4=4$, $\abs{\vec a}=\abs{\vec b}=3$, takže
$\cos\varphi=\dfrac{4}{9}\doteq0\dc444$, $\varphi\doteq63\dc6^\circ$.
\end{priklad}

\begin{priklad}{Kolmost s~parametrem}
Pro jakou hodnotu $k$ jsou vektory $\vec a=(1,k,2)$ a $\vec b=(3,-1,1)$ kolmé?

\medskip\noindent\textbf{Řešení.} $\vec a\cdot\vec b=3-k+2=0\ \Rightarrow\ k=5$.
\end{priklad}

\begin{priklad}{Průmět a kolmá složka}
$\vec a=(4,3)$, $\vec b=(2,-1)$. Určete průmět $\vec a$ do směru $\vec b$ a kolmou složku.

\medskip\noindent\textbf{Řešení.}
$\vec a\cdot\vec b=8-3=5$, $\vec b\cdot\vec b=5$, takže
\begin{align*}
  \proj_{\vec b}\vec a&=\frac55\,(2,-1)=(2,-1),\\
  a_{\parallel b}&=\frac{5}{\sqrt5}=\sqrt5\doteq2\dc24,\\
  \vec a_\perp&=(4,3)-(2,-1)=(2,4).
\end{align*}
Kontrola: $\vec a_\perp\cdot\vec b=4-4=0$. \checkmark
\end{priklad}

\begin{priklad}{Práce síly}
Síla $\vec F=(10,5,2)\jed{N}$ posune těleso po dráze $\vec s=(4,3,1)\jed{m}$.
Práce: $W=\vec F\cdot\vec s=40+15+2=57\jed{J}$.
\end{priklad}

% =====================================================================
\section{Vektorový součin}
% =====================================================================

Vektorový součin existuje (v~této podobě) jen v~\emph{trojrozměrném} prostoru.
Na rozdíl od skalárního součinu je jeho výsledkem \textbf{vektor}.

\subsection{Definice a geometrický význam}

\begin{definice}{Vektorový součin}
\textbf{Vektorový součin} vektorů $\vec a,\vec b$ je vektor $\vec c=\vec a\times\vec b$ s~vlastnostmi:
\begin{enumerate}[label=(\arabic*),nosep]
  \item \textbf{velikost:} $\abs{\vec a\times\vec b}=\abs{\vec a}\,\abs{\vec b}\sin\varphi$,
        kde $\varphi\in\langle0^\circ,180^\circ\rangle$ je úhel vektorů;
  \item \textbf{směr:} $\vec c$ je kolmý na $\vec a$ i~na $\vec b$;
  \item \textbf{orientace:} trojice $(\vec a,\vec b,\vec c)$ tvoří \emph{pravotočivou} soustavu
        (pravidlo pravé ruky, viz níže).
\end{enumerate}
Jsou-li $\vec a,\vec b$ kolineární nebo je jeden z~nich nulový, je $\vec a\times\vec b=\vec o$.
\end{definice}

\textbf{Pravidlo pravé ruky:} prsty pravé ruky ohněte od prvního vektoru $\vec a$
k~druhému $\vec b$ (kratší cestou); vztyčený palec ukazuje směr $\vec a\times\vec b$.
Ekvivalentně: otáčíme-li pravotočivým šroubem od $\vec a$ k~$\vec b$, šroub se pohybuje
ve směru $\vec a\times\vec b$.

\begin{figure}[!htbp]
\centering
\begin{tikzpicture}[x={(-0.6cm,-0.45cm)},y={(1.2cm,0cm)},z={(0cm,1.2cm)}]
  \fill[veccol1!12] (0,0,0) -- (3,0.5,0) -- (3.5,3,0) -- (0.5,2.5,0) -- cycle;
  \draw[helpl] (3,0.5,0) -- (3.5,3,0) -- (0.5,2.5,0);
  \draw[vec]  (0,0,0) -- (3,0.5,0) node[below left]{$\vec a$};
  \draw[vecg] (0,0,0) -- (0.5,2.5,0) node[above right]{$\vec b$};
  \draw[vecr] (0,0,0) -- (0,0,3.2) node[above]{$\vec a\times\vec b$};
  \node[font=\small] at (1.8,1.5,0) {$S=\abs{\vec a\times\vec b}$};
  \node[font=\small] at (0.55,0.45,0) {$\varphi$};
\end{tikzpicture}
\caption{Vektor $\vec a\times\vec b$ je kolmý na rovinu vektorů $\vec a,\vec b$;
jeho velikost je obsah rovnoběžníku, který vektory určují.}
\end{figure}

\subsection{Výpočet pomocí souřadnic}

Z~vlastností báze $\ii\times\jj=\kk$, $\jj\times\kk=\ii$, $\kk\times\ii=\jj$
(a~při opačném pořadí opačné znaménko) a $\ii\times\ii=\vec o$ atd.\ plyne
\begin{equation}
  \vec a\times\vec b=
  \begin{vmatrix}
    \ii & \jj & \kk\\
    a_x & a_y & a_z\\
    b_x & b_y & b_z
  \end{vmatrix}
  =\bigl(a_yb_z-a_zb_y,\ \ a_zb_x-a_xb_z,\ \ a_xb_y-a_yb_x\bigr).
\end{equation}
Determinant rozvíjíme podle prvního řádku (pozor na \emph{mínus} u~složky $\jj$!).

\begin{poznamka}
Mnemotechnika pro bázi: $\ii\to\jj\to\kk\to\ii$ jsou „cyklicky vpřed“: každá
dvojice v~pořadí dává třetí se znaménkem $+$, v~opačném pořadí se znaménkem $-$.
Například $\jj\times\ii=-\kk$.
\end{poznamka}

\subsection{Vlastnosti}

\begin{veta}{Vlastnosti vektorového součinu}
\begin{enumerate}[label=(\roman*),nosep]
  \item $\vec a\times\vec b=-(\vec b\times\vec a)$ \hfill \textbf{(antikomutativita)}
  \item $\vec a\times\vec a=\vec o$
  \item $(k\vec a)\times\vec b=k(\vec a\times\vec b)=\vec a\times(k\vec b)$
  \item $\vec a\times(\vec b+\vec c)=\vec a\times\vec b+\vec a\times\vec c$ \hfill (distributivita)
  \item $\vec a\times\vec b=\vec o\iff\vec a,\vec b$ jsou kolineární (nenulové vektory)
  \item $(\vec a\times\vec b)\cdot\vec a=0$ a $(\vec a\times\vec b)\cdot\vec b=0$
\end{enumerate}
\end{veta}

\begin{pozor}
Vektorový součin je \textbf{nekomutativní} a \textbf{neasociativní}:
obecně $(\vec a\times\vec b)\times\vec c\neq\vec a\times(\vec b\times\vec c)$.
Například $(\ii\times\ii)\times\jj=\vec o$, ale $\ii\times(\ii\times\jj)=\ii\times\kk=-\jj$.
Pořadí a závorky proto \emph{vždy} hlídejte.
\end{pozor}

\subsection{Geometrické využití}

\begin{align}
  S_{\text{rovnoběžníku}}&=\abs{\vec a\times\vec b},\\
  S_{\text{trojúhelníku }ABC}&=\tfrac12\abs{\overrightarrow{AB}\times\overrightarrow{AC}},\\
  \vec n&=\vec a\times\vec b\quad\text{je normálový vektor roviny určené }\vec a,\vec b.
\end{align}

\begin{poznamka}
Ve 2D ($a_z=b_z=0$) má vektorový součin jedinou nenulovou složku:
$\vec a\times\vec b=(0,0,\,a_xb_y-a_yb_x)$. Číslo $a_xb_y-a_yb_x$ je \emph{orientovaný
obsah} rovnoběžníku: je kladné, když se od $\vec a$ k~$\vec b$ otáčíme proti směru
hodinových ručiček.
\end{poznamka}

\begin{ekt}[Vektorový součin ve fyzice a elektrotechnice]
\begin{itemize}
  \item \textbf{Moment síly} k~bodu: $\vec M=\vec r\times\vec F$.
  \item \textbf{Lorentzova síla} na náboj $q$ v~magnetickém poli: $\vec F=q\,\vec v\times\vec B$.
        (Se silou v~elektrickém poli dohromady: $\vec F=q(\vec E+\vec v\times\vec B)$.)
  \item \textbf{Síla na vodič} s~proudem: $\vec F=I\,\vec l\times\vec B$.
  \item \textbf{Poyntingův vektor} (hustota toku výkonu elmag.\ vlny):
        $\vec S=\vec E\times\vec H$. Ukazuje směr šíření vlny.
  \item \textbf{Obvodová rychlost} rotujícího tělesa: $\vec v=\vec\omega\times\vec r$.
  \item \textbf{Biotův--Savartův zákon}: $\mathrm d\vec B=\frac{\mu_0 I}{4\pi}\,\frac{\mathrm d\vec l\times\vec r}{r^3}$.
\end{itemize}
\end{ekt}

\begin{priklad}{Výpočet vektorového součinu}
$\vec a=(1,2,3)$, $\vec b=(4,5,6)$.

\medskip\noindent\textbf{Řešení.}
\[
  \vec a\times\vec b=
  \begin{vmatrix}\ii&\jj&\kk\\1&2&3\\4&5&6\end{vmatrix}
  =(2\cdot6-3\cdot5,\ 3\cdot4-1\cdot6,\ 1\cdot5-2\cdot4)=(-3,\,6,\,-3).
\]
Kontrola kolmosti: $\vec a\cdot(\vec a\times\vec b)=-3+12-9=0$ a
$\vec b\cdot(\vec a\times\vec b)=-12+30-18=0$. \checkmark\quad
Velikost: $\abs{\vec a\times\vec b}=\sqrt{9+36+9}=\sqrt{54}=3\sqrt6\doteq7\dc35$.
\end{priklad}

\begin{priklad}{Obsah trojúhelníku a rovina}
Trojúhelník má vrcholy $A=[1,0,0]$, $B=[0,2,0]$, $C=[0,0,3]$. Určete jeho obsah
a normálový vektor roviny.

\medskip\noindent\textbf{Řešení.} $\overrightarrow{AB}=(-1,2,0)$, $\overrightarrow{AC}=(-1,0,3)$,
\[
  \overrightarrow{AB}\times\overrightarrow{AC}
  =(2\cdot3-0\cdot0,\ 0\cdot(-1)-(-1)\cdot3,\ (-1)\cdot0-2\cdot(-1))=(6,3,2).
\]
$\abs{(6,3,2)}=\sqrt{36+9+4}=7$, proto $S=\tfrac12\cdot7=3\dc5$.
Normálový vektor: $\vec n=(6,3,2)$; rovina má rovnici $6x+3y+2z=6$
(dosazením bodů $A,B,C$ ověříme). \checkmark
\end{priklad}

\begin{priklad}{Lorentzova síla}
Proton ($q=1\dc6\cdot10^{-19}\jed{C}$) letí rychlostí $\vec v=(2\cdot10^{6},0,0)\jed{m/s}$
v~poli $\vec B=(0,0,0\dc5)\jed{T}$. Určete sílu.

\medskip\noindent\textbf{Řešení.}
$\vec v\times\vec B=\bigl(0\cdot0{,}5-0\cdot0,\ 0\cdot0-2\cdot10^6\cdot0{,}5,\ 0\bigr)=(0,\,-10^6,\,0)$, tedy
\[
  \vec F=q\,\vec v\times\vec B=(0,\,-1\dc6\cdot10^{-13},\,0)\jed{N}.
\]
Síla je kolmá na rychlost i~na pole, proto proton opisuje kružnici.
\end{priklad}

\begin{priklad}{Moment síly}
Síla $\vec F=(0,0,10)\jed{N}$ působí v~bodě $\vec r=(0,2,0)\jed{m}$.
$\vec M=\vec r\times\vec F=(2\cdot10-0,\ 0-0,\ 0)=(20,0,0)\jed{N\cdot m}$.
\end{priklad}

% =====================================================================
\section{Smíšený a dvojitý vektorový součin}
% =====================================================================

\subsection{Smíšený součin}\label{sec:smiseny}

\begin{definice}{Smíšený součin}
\textbf{Smíšený součin} tří vektorů je číslo
\begin{equation}
  (\vec a\times\vec b)\cdot\vec c=
  \begin{vmatrix}
    a_x&a_y&a_z\\ b_x&b_y&b_z\\ c_x&c_y&c_z
  \end{vmatrix}.
\end{equation}
\end{definice}

\begin{figure}[!htbp]
\centering
\begin{tikzpicture}[x={(-0.6cm,-0.45cm)},y={(1.2cm,0cm)},z={(0cm,1.2cm)}]
  \coordinate (O)   at (0,0,0);
  \coordinate (A)   at (2,0.3,0);
  \coordinate (B)   at (0.3,2.5,0);
  \coordinate (C)   at (0.4,0.4,2);
  \coordinate (AB)  at (2.3,2.8,0);
  \coordinate (AC)  at (2.4,0.7,2);
  \coordinate (BC)  at (0.7,2.9,2);
  \coordinate (ABC) at (2.7,3.2,2);
  \fill[veccol2!15] (C) -- (AC) -- (ABC) -- (BC) -- cycle;
  \fill[veccol1!10] (A) -- (AB) -- (ABC) -- (AC) -- cycle;
  \draw[gray!70] (O)--(A)--(AB)--(B)--cycle;
  \draw[gray!70] (C)--(AC)--(ABC)--(BC)--cycle;
  \draw[gray!70] (O)--(C) (A)--(AC) (B)--(BC) (AB)--(ABC);
  \draw[vec]  (O) -- (A) node[below left]{$\vec a$};
  \draw[vecg] (O) -- (B) node[below right]{$\vec b$};
  \draw[vecr] (O) -- (C) node[above left]{$\vec c$};
\end{tikzpicture}
\caption{Rovnoběžnostěn určený vektory $\vec a,\vec b,\vec c$.}
\end{figure}

\begin{veta}{Vlastnosti a význam}
\begin{itemize}
  \item \textbf{Objem rovnoběžnostěnu:} $V=\abs{(\vec a\times\vec b)\cdot\vec c}$;
        \textbf{objem čtyřstěnu} (jehlanu s~trojúhelníkovou podstavou):
        $V=\tfrac16\abs{(\vec a\times\vec b)\cdot\vec c}$.
  \item \textbf{Cyklická záměna} nemění hodnotu:
        $(\vec a\times\vec b)\cdot\vec c=(\vec b\times\vec c)\cdot\vec a=(\vec c\times\vec a)\cdot\vec b$;
        prohození dvou vektorů změní znaménko. Také $(\vec a\times\vec b)\cdot\vec c=\vec a\cdot(\vec b\times\vec c)$.
  \item \textbf{Komplanarita:} vektory $\vec a,\vec b,\vec c$ jsou komplanární
        (lineárně závislé) $\iff(\vec a\times\vec b)\cdot\vec c=0$.
  \item Kladný smíšený součin značí, že $(\vec a,\vec b,\vec c)$ tvoří pravotočivou soustavu.
\end{itemize}
\end{veta}

\begin{priklad}{Objem a komplanarita}
\textbf{a)} $\vec a=(1,0,0)$, $\vec b=(1,2,0)$, $\vec c=(1,1,3)$:
\[
  (\vec a\times\vec b)\cdot\vec c=
  \begin{vmatrix}1&0&0\\1&2&0\\1&1&3\end{vmatrix}=1\cdot(2\cdot3-0\cdot1)=6,
\]
tedy $V_{\text{rovnoběžnostěnu}}=6$ a $V_{\text{čtyřstěnu}}=1$.

\smallskip\noindent
\textbf{b)} Jsou vektory $\vec u=(1,2,3)$, $\vec v=(4,5,6)$, $\vec w=(7,8,9)$ komplanární?
\[
  \begin{vmatrix}1&2&3\\4&5&6\\7&8&9\end{vmatrix}
  =1(45-48)-2(36-42)+3(32-35)=-3+12-9=0 .
\]
Ano, leží v~jedné rovině (skutečně $\vec w=2\vec v-\vec u$).
\end{priklad}

\subsection{Dvojitý vektorový součin}

\begin{veta}{„BAC--CAB“}
$\vec a\times(\vec b\times\vec c)=\vec b\,(\vec a\cdot\vec c)-\vec c\,(\vec a\cdot\vec b)$
\end{veta}

Vzorec se hodí při úpravách výrazů (např.\ v~elektrodynamice). Zapamatujete si
ho jako „prostřední krát skalární součin zbylých, minus poslední krát skalární
součin zbylých“. Dále platí \textbf{Lagrangeova identita}
\[
  \abs{\vec a\times\vec b}^2+(\vec a\cdot\vec b)^2=\abs{\vec a}^2\abs{\vec b}^2,
\]
která je jen zápisem vztahu $\sin^2\varphi+\cos^2\varphi=1$.

% =====================================================================
\section{Skládání a rozklad vektorů}\label{sec:skladani}
% =====================================================================

\subsection{Výslednice}

Působí-li na těleso několik sil $\vec F_1,\dots,\vec F_n$, nahradíme je jedinou silou
(\textbf{výslednicí}) $\vec F$, která má stejný účinek:
\begin{equation}
  \vec F=\sum_{i=1}^{n}\vec F_i,\qquad
  F_x=\sum F_{ix},\quad F_y=\sum F_{iy},\quad F_z=\sum F_{iz}.
\end{equation}
Skládání vektorů je tedy prostě jejich sčítání. Prakticky postupujeme takto:
\begin{enumerate}
  \item Zvolíme vhodnou soustavu souřadnic.
  \item \textbf{Rozložíme} každý vektor na složky:
        $F_{ix}=F_i\cos\varphi_i,\ F_{iy}=F_i\sin\varphi_i$.
  \item Složky po osách \textbf{sečteme}.
  \item Ze složek výslednice určíme velikost
        $F=\sqrt{F_x^2+F_y^2}$ a směr $\varphi=\operatorname{atan2}(F_y,F_x)$.
\end{enumerate}

Pro \emph{dva} vektory lze velikost součtu získat přímo z~definice skalárního součinu
(kosinová věta):
\begin{equation}
  \abs{\vec a+\vec b}^2=(\vec a+\vec b)\cdot(\vec a+\vec b)=a^2+b^2+2ab\cos\varphi,
  \qquad
  \abs{\vec a-\vec b}^2=a^2+b^2-2ab\cos\varphi,
\end{equation}
kde $\varphi$ je úhel, který vektory svírají při společném počátku.

\begin{priklad}{Dvě síly}
Síly $F_1=30\jed{N}$ a~$F_2=40\jed{N}$ svírají úhel $60^\circ$. Velikost výslednice:
\[
  F=\sqrt{30^2+40^2+2\cdot30\cdot40\cos60^\circ}=\sqrt{900+1600+1200}=\sqrt{3700}\doteq60\dc8\jed{N}.
\]
Pro $\varphi=90^\circ$ by vyšlo $F=50\jed{N}$ (Pythagorova věta), pro $\varphi=0^\circ$
je $F=70\jed{N}$ (maximum), pro $\varphi=180^\circ$ je $F=10\jed{N}$ (minimum).
\end{priklad}

\begin{priklad}{Tři síly v~rovině}
Na bod působí síly: $F_1=50\jed{N}$ ve směru $0^\circ$, $F_2=80\jed{N}$ ve směru
$60^\circ$, $F_3=30\jed{N}$ ve směru $150^\circ$ (úhly od osy $x$). Určete výslednici.

\medskip\noindent\textbf{Řešení.}
\begin{align*}
  \vec F_1&=(50,\ 0),\\
  \vec F_2&=(80\cos60^\circ,\ 80\sin60^\circ)=(40,\ 69\dc28),\\
  \vec F_3&=(30\cos150^\circ,\ 30\sin150^\circ)=(-25\dc98,\ 15),\\
  \vec F&=\vec F_1+\vec F_2+\vec F_3=(64\dc02,\ 84\dc28)\jed{N}.
\end{align*}
Velikost $F=\sqrt{64\dc02^2+84\dc28^2}\doteq105\dc8\jed{N}$, směr
$\varphi=\operatorname{arctg}\dfrac{84\dc28}{64\dc02}\doteq52\dc8^\circ$.
\end{priklad}

\subsection{Rozklad vektoru na složky}

Opačná úloha: daný vektor rozložit do zadaných směrů.

\begin{itemize}
  \item \textbf{Do kolmých směrů} (např.\ osy $x,y$ nebo směr rovnoběžný a kolmý
        k~podložce): složky získáme skalárním součinem s~jednotkovými vektory
        směrů, $F_{\parallel}=\vec F\cdot\vec e_{\parallel}$, $F_{\perp}=\vec F\cdot\vec e_{\perp}$
        (viz průmět v~kapitole o~skalárním součinu).
  \item \textbf{Do dvou obecných směrů} $\vec u,\vec v$: hledáme $\alpha,\beta$
        z~rovnice $\vec F=\alpha\vec u+\beta\vec v$, což vede na soustavu dvou rovnic
        (viz Příklad o~rozkladu do báze), nebo lze užít sinovou větu v~trojúhelníku sil.
\end{itemize}

\begin{figure}[!htbp]
\centering
\begin{tikzpicture}[scale=1.2]
  \coordinate (A) at (0,0);
  \coordinate (B) at (5,0);
  \coordinate (C) at (5,2.887);
  \draw[thick,fill=gray!15] (A) -- (B) -- (C) -- cycle;
  \coordinate (M) at (2.5,1.4434);
  \begin{scope}[shift={(M)},rotate=30]
    \draw[thick,fill=blue!15] (-0.5,0) rectangle (0.5,0.7);
  \end{scope}
  \coordinate (S) at ($(M)+(-0.175,0.303)$);
  \draw[helpl] ($(S)+(-0.866,-0.5)$) -- ($(S)+(0,-2)$) -- ($(S)+(0.866,-1.5)$);
  \draw[vecr] (S) -- ++(0,-2)         node[below]{$\vec G$};
  \draw[vec]  (S) -- ++(-0.866,-0.5)  node[below left]{$\vec G_t$};
  \draw[vecg] (S) -- ++(0.866,-1.5)   node[below right]{$\vec G_n$};
  \draw pic[draw,pic text={\(\alpha\)},angle radius=1cm,angle eccentricity=1.4] {angle=B--A--C};
\end{tikzpicture}
\caption{Rozklad tíhové síly $\vec G$ na nakloněné rovině do složky $\vec G_t$
rovnoběžné s~rovinou a složky $\vec G_n$ kolmé k~rovině. Platí
$G_t=G\sin\alpha$, $G_n=G\cos\alpha$.}
\end{figure}

\begin{priklad}{Nakloněná rovina}
Těleso o~hmotnosti $m=10\jed{kg}$ leží na nakloněné rovině se sklonem $\alpha=30^\circ$
($g=9\dc81\jed{m/s^2}$). Určete složky tíhy.

\medskip\noindent\textbf{Řešení.} $G=mg=98\dc1\jed{N}$,
\[
  G_t=G\sin\alpha=98\dc1\cdot0\dc5\doteq49\dc05\jed{N},\qquad
  G_n=G\cos\alpha=98\dc1\cdot0\dc866\doteq84\dc96\jed{N}.
\]
Kontrola: $\sqrt{49\dc05^2+84\dc96^2}\doteq98\dc1\jed{N}$. \checkmark
\end{priklad}

\subsection{Rovnováha sil}

Těleso (hmotný bod) je v~rovnováze, je-li \textbf{výslednice sil nulová}:
$\sum\vec F_i=\vec o$. Ve složkách to jsou 2 rovnice (v~rovině) resp.\ 3 rovnice
(v~prostoru): $\sum F_{ix}=0,\ \sum F_{iy}=0,\ \sum F_{iz}=0$. Znamená to, že vektory sil
tvoří po zřetězení \emph{uzavřený mnohoúhelník}.

\begin{priklad}{Závaží na dvou lanech}
Závaží o~tíze $G=100\jed{N}$ visí na dvou stejně dlouhých lanech, každé je odkloněno
o~$30^\circ$ od vodorovné roviny. Určete tahové síly v~lanech.

\medskip\noindent\textbf{Řešení.} Ze symetrie jsou tahy stejné, $T_1=T_2=T$.
\[
  \vec T_1=(-T\cos30^\circ,\ T\sin30^\circ),\quad
  \vec T_2=(T\cos30^\circ,\ T\sin30^\circ),\quad
  \vec G=(0,\,-100).
\]
Podmínka rovnováhy $\vec T_1+\vec T_2+\vec G=\vec o$: složka $x$ je splněna automaticky,
složka $y$ dává $2T\sin30^\circ-100=0\Rightarrow T=100\jed{N}$.
\end{priklad}

\subsection{Skládání fázorů ve střídavých obvodech}

\begin{ekt}[Fázory: vektory rotující v~rovině]
Harmonický signál $u(t)=U_m\cos(\omega t+\varphi)$ lze znázornit \textbf{fázorem}:
vektorem (komplexním číslem) o~velikosti $U_m$ a úhlu $\varphi$ v~rovině. Všechny
signály téhož kmitočtu se „točí“ stejnou rychlostí, takže jejich vzájemná poloha
zůstává a stačí pracovat s~pevnými vektory. Součet signálů (například napětí na
sériově zapojených prvcích) odpovídá \textbf{vektorovému součtu fázorů}.
\end{ekt}

\begin{figure}[!htbp]
\centering
\begin{tikzpicture}[scale=0.85]
  \draw[ax] (-0.5,0) -- (4.8,0) node[right]{Re};
  \draw[ax] (0,-0.5) -- (0,5.2) node[above]{Im\,($\jimag$)};
  \draw[vec]  (0,0) -- (3,0)   node[below]{$\hat U_1=3$};
  \draw[vecg] (0,0) -- (0,4)   node[left]{$\hat U_2=4\jimag$};
  \draw[helpl] (3,0) -- (3,4) -- (0,4);
  \draw[vecr] (0,0) -- (3,4)   node[midway,above left]{$\hat U=5\,\angle53\dc13^\circ$};
  \coordinate (O) at (0,0);
  \coordinate (R) at (1,0);
  \coordinate (U) at (3,4);
  \draw pic[draw,pic text={\(\varphi\)},angle radius=8mm,angle eccentricity=1.5] {angle=R--O--U};
\end{tikzpicture}
\caption{Fázorový diagram: $u_1=3\cos\omega t$ a $u_2=4\cos(\omega t+90^\circ)$.}
\end{figure}

\begin{priklad}{Součet dvou napětí}
V~sériovém obvodu je $u_1(t)=3\cos\omega t$ [V] a $u_2(t)=4\cos(\omega t+90^\circ)$ [V].
Určete celkové napětí $u(t)$.

\medskip\noindent\textbf{Řešení.} Fázory: $\hat U_1=3$, $\hat U_2=4\,\mathrm e^{\jimag90^\circ}=4\jimag$.
Součet
\[
  \hat U=3+4\jimag,\qquad \abs{\hat U}=\sqrt{3^2+4^2}=5,\qquad
  \varphi=\operatorname{arctg}\tfrac43\doteq53\dc13^\circ .
\]
Tedy $u(t)=5\cos(\omega t+53\dc13^\circ)$ V. Všimněte si, že amplitudy se nesčítají
algebraicky ($3+4\neq5$), protože signály nejsou ve fázi.
\end{priklad}

% =====================================================================
\section{Souhrn vzorců a časté chyby}
% =====================================================================

\begin{tcolorbox}[enhanced,breakable,colback=primary!4,colframe=primary,
  title={\bfseries Přehled vzorců (3D)},fonttitle=\bfseries,coltitle=white,colbacktitle=primary]
\textbf{Velikost, jednotkový vektor, vektor mezi body}
\[
  \abs{\vec a}=\sqrt{a_x^2+a_y^2+a_z^2},\quad
  \vec a^{\,0}=\frac{\vec a}{\abs{\vec a}},\quad
  \overrightarrow{AB}=B-A
\]
\textbf{Lineární operace (po složkách)}
\[
  \vec a\pm\vec b=(a_x\pm b_x,\,a_y\pm b_y,\,a_z\pm b_z),\qquad
  k\vec a=(ka_x,\,ka_y,\,ka_z)
\]
\textbf{Skalární součin} (výsledek je \emph{číslo})
\[
  \vec a\cdot\vec b=a_xb_x+a_yb_y+a_zb_z=\abs{\vec a}\abs{\vec b}\cos\varphi,\qquad
  \vec a\perp\vec b\iff\vec a\cdot\vec b=0
\]
\textbf{Vektorový součin} (výsledek je \emph{vektor})
\[
  \vec a\times\vec b=(a_yb_z-a_zb_y,\ a_zb_x-a_xb_z,\ a_xb_y-a_yb_x),\qquad
  \abs{\vec a\times\vec b}=\abs{\vec a}\abs{\vec b}\sin\varphi
\]
\textbf{Smíšený součin} (výsledek je \emph{číslo})
\[
  (\vec a\times\vec b)\cdot\vec c=\det\begin{pmatrix}\vec a\\ \vec b\\ \vec c\end{pmatrix},\qquad
  V=\abs{(\vec a\times\vec b)\cdot\vec c}
\]
\textbf{Použití}
\[
  \cos\varphi=\frac{\vec a\cdot\vec b}{\abs{\vec a}\abs{\vec b}},\quad
  \proj_{\vec b}\vec a=\frac{\vec a\cdot\vec b}{\abs{\vec b}^2}\vec b,\quad
  S_\triangle=\tfrac12\abs{\vec a\times\vec b}
\]
\end{tcolorbox}

\subsection{Přehled operací}

\begin{center}
\small
\begin{tabular}{@{}llll@{}}
\toprule
\textbf{Operace} & \textbf{Značení} & \textbf{Výsledek} & \textbf{Komutativní?}\\
\midrule
Sčítání            & $\vec a+\vec b$            & vektor & ano\\
Násobení skalárem  & $k\vec a$                  & vektor & ano ($k\vec a=\vec a k$)\\
Skalární součin    & $\vec a\cdot\vec b$        & \textbf{skalár} & ano\\
Vektorový součin   & $\vec a\times\vec b$       & \textbf{vektor} & \textbf{ne} (mění znaménko)\\
Smíšený součin     & $(\vec a\times\vec b)\cdot\vec c$ & skalár & cyklicky ano\\
\bottomrule
\end{tabular}
\end{center}

\subsection{Nejčastější chyby}

\begin{enumerate}
  \item Psaní $\abs{\vec a+\vec b}=\abs{\vec a}+\abs{\vec b}$ (platí jen pro souhlasně rovnoběžné vektory).
  \item Záměna \emph{skalárního} a \emph{vektorového} součinu, případně tečka místo
        křížku (výsledek je jednou číslo, jednou vektor).
  \item Zapomenuté \emph{mínus} u~prostřední složky determinantu při výpočtu $\vec a\times\vec b$.
  \item Nesprávný kvadrant úhlu při použití $\operatorname{arctg}$.
  \item Odčítání „začátek minus konec“ místo \textbf{konec minus začátek} u~$\overrightarrow{AB}$.
  \item Zaměnění pořadí ve vektorovém součinu (a~tedy znaménka výsledku).
  \item Kalkulačka v~nesprávném režimu (DEG/RAD) při počítání s~úhly.
  \item Dělení vektoru vektorem nebo psaní $\vec a>\vec b$: takové operace neexistují.
\end{enumerate}

% =====================================================================
\section{Příklady k~procvičení}
% =====================================================================

\begin{enumerate}[label=\textbf{\arabic*.},leftmargin=*,itemsep=3pt]
  \item Body $A=[3,-1,2]$ a $B=[6,3,-2]$. Určete $\overrightarrow{AB}$, $\abs{\overrightarrow{AB}}$
        a jednotkový vektor ve směru $\overrightarrow{AB}$.
  \item Pro $\vec a=(2,-1,3)$ a $\vec b=(1,4,-2)$ vypočtěte
        $\vec a+\vec b$, $\ \vec a-\vec b$, $\ 3\vec a-2\vec b$, $\ \vec a\cdot\vec b$, $\ \vec a\times\vec b$
        a úhel mezi vektory.
  \item Najděte jednotkový vektor ve směru $(6,-2,3)$.
  \item Pro které $k$ jsou vektory $(k,2,-1)$ a $(3,-1,4)$ kolmé?
  \item Dokažte, že vektory $(1,2,-1)$ a $(-3,-6,3)$ jsou kolineární.
  \item Rozložte $\vec a=(7,1)$ na složku rovnoběžnou a složku kolmou k~$\vec b=(3,4)$.
  \item Určete obsah trojúhelníku $A=[0,0,0]$, $B=[2,1,0]$, $C=[1,3,0]$.
  \item Určete objem čtyřstěnu s~vrcholy $A=[0,0,0]$, $B=[1,0,0]$, $C=[0,2,0]$, $D=[0,0,3]$.
  \item Síly $\vec F_1=(3,4,0)$, $\vec F_2=(-1,2,5)$, $\vec F_3=(0,-6,-5)$ (v newtonech)
        působí v~jednom bodě. Určete výslednici a její velikost.
  \item Síla $\vec F=(0,0,4)\jed{N}$ působí v~bodě $\vec r=(1,2,0)\jed{m}$. Určete moment síly
        vzhledem k~počátku.
  \item Napětí $u_1(t)=6\cos\omega t$ a $u_2(t)=8\cos(\omega t-90^\circ)$ [V] se sčítají.
        Určete amplitudu a fázi výsledného napětí.
  \item Je dáno $\abs{\vec a}=4$, $\abs{\vec b}=3$, úhel $\varphi=60^\circ$. Vypočtěte
        $\vec a\cdot\vec b$, $\abs{\vec a\times\vec b}$ a $\abs{\vec a+\vec b}$.
\end{enumerate}

\subsection*{Výsledky}

\begin{enumerate}[label=\textbf{\arabic*.},leftmargin=*,itemsep=1pt]
  \item $\overrightarrow{AB}=(3,4,-4)$, $\ \abs{\overrightarrow{AB}}=\sqrt{41}\doteq6\dc40$,
        $\ \overrightarrow{AB}^{\,0}=\frac1{\sqrt{41}}(3,4,-4)$.
  \item $\vec a+\vec b=(3,3,1)$, $\ \vec a-\vec b=(1,-5,5)$, $\ 3\vec a-2\vec b=(4,-11,13)$,
        $\ \vec a\cdot\vec b=-8$, $\ \vec a\times\vec b=(-10,7,9)$;
        $\cos\varphi=\dfrac{-8}{\sqrt{14}\sqrt{21}}\doteq-0\dc467$, $\varphi\doteq117\dc8^\circ$.
  \item $\left(\tfrac67,-\tfrac27,\tfrac37\right)$.
  \item $3k-2-4=0\Rightarrow k=2$.
  \item $(-3,-6,3)=-3\cdot(1,2,-1)$.
  \item $\vec a\cdot\vec b=25=\abs{\vec b}^2$, takže $\vec a_\parallel=(3,4)$, $\ \vec a_\perp=(4,-3)$.
  \item $\overrightarrow{AB}\times\overrightarrow{AC}=(0,0,5)$, $S=\tfrac52=2\dc5$.
  \item $V=\tfrac16\abs{(\overrightarrow{AB}\times\overrightarrow{AC})\cdot\overrightarrow{AD}}=\tfrac16\cdot6=1$.
  \item $\vec F=(2,0,0)\jed{N}$, $\ F=2\jed{N}$.
  \item $\vec M=\vec r\times\vec F=(8,-4,0)\jed{N\cdot m}$.
  \item $\hat U=6-8\jimag$, $\ U_m=10\jed{V}$, $\ \varphi=-53\dc13^\circ$.
  \item $\vec a\cdot\vec b=6$, $\ \abs{\vec a\times\vec b}=6\sqrt3\doteq10\dc39$,
        $\ \abs{\vec a+\vec b}=\sqrt{16+9+12}=\sqrt{37}\doteq6\dc08$.
\end{enumerate}

\vspace{1.2cm}
\begin{center}
\color{primary}\rule{0.4\textwidth}{0.6pt}\\[4pt]
\small\itshape Hodně štěstí při studiu!
\end{center}

\end{document}